% GENERATED FILE. Edit canonical lessons, then run npm run generate:books.
% canonical-source-hash: fnv1a64:775b698030d2a937
% canonical-lessons: GE-C503-gaertnern, GE-C503-unkraut, GE-C503-jaeten, GE-C503-giesskanne, GE-C503-kompost

\chapter{Im Garten}
\label{ch:ge-im-garten}
\hlchaptermodality{503}

\begin{chapteropening}
\textbf{By the end of this chapter:} \emph{I can talk about gardening.}

Complete the last lesson of chapter 503: I can talk about gardening.
\end{chapteropening}

\section[gärtnern]{gärtnern — to garden}

\label{lesson:GE-C503-gaertnern}



\begin{quote}
Before the new one: say the German for lens, then the German for to snap.
\end{quote}

\subsection*{You'll want to know: gärtnern}
\textbf{gärtnern} — \textquotedblleft{}to garden\textquotedblright{}.

\begin{quote}
In ihrer Freizeit gärtnert meine Mutter gern.
\end{quote}

\begin{quote}
\emph{In her free time my mother likes gardening.}
\end{quote}

Built on \emph{der Gärtner} — gardener.

\begin{cousinweb}
\emph{gärtnern} is built on \emph{Gärtner} \textquotedblleft{}gardener\textquotedblright{}, from \emph{Garten}, which is inherited from Old High German \emph{garto} \textquotedblleft{}garden, enclosure\textquotedblright{}, from Proto-Germanic \emph{gardô} (EtymDB).
\end{cousinweb}

\begin{grammarlens}[title={What you've built}]
One more word for telling, explaining and describing at B1.
\end{grammarlens}

\subsection*{Your turn}
\begin{itemize}
\raggedright
  \item \emph{Say it:} \emph{gärtnern}
  \item \emph{Say it:} \emph{gärtnern}, once more
  \item \emph{Say it:} a sentence of your own with \emph{gärtnern}
  \item \emph{Recall:} say the German for lens, then the German for to snap, then say \emph{gärtnern} again
\end{itemize}

\begin{tcolorbox}[breakable,colback=teal!4,colframe=teal!35!black,title={Before you move on}]
What does gärtnern mean? (\textquotedblleft{}To garden\textquotedblright{}.) Say it once more.
\end{tcolorbox}
\section[das Unkraut]{das Unkraut — weeds}

\label{lesson:GE-C503-unkraut}



\begin{quote}
Before the new one: say the German for to snap, then the German for to garden.
\end{quote}

\subsection*{You'll want to know: das Unkraut}
\textbf{das Unkraut} — \textquotedblleft{}weeds\textquotedblright{}.

\begin{quote}
Zwischen den Tomaten wächst viel Unkraut.
\end{quote}

\begin{quote}
\emph{Lots of weeds are growing between the tomatoes.}
\end{quote}

Mostly singular for all the weeds in a bed. Plural: \emph{die Unkräuter}.

\begin{cousinweb}
\emph{Unkraut} is built on \emph{un-} \textquotedblleft{}not, bad\textquotedblright{} and \emph{Kraut} \textquotedblleft{}herb, plant\textquotedblright{}, inherited from Old High German \emph{krūt}, from Proto-Germanic \emph{krūdą} (EtymDB).
\end{cousinweb}

\begin{grammarlens}[title={What you've built}]
One more word for telling, explaining and describing at B1.
\end{grammarlens}

\subsection*{Your turn}
\begin{itemize}
\raggedright
  \item \emph{Say it:} \emph{das Unkraut}
  \item \emph{Say it:} \emph{das Unkraut}, once more
  \item \emph{Say it:} a sentence of your own with \emph{das Unkraut}
  \item \emph{Recall:} say the German for to snap, then the German for to garden, then say \emph{das Unkraut} again
\end{itemize}

\begin{tcolorbox}[breakable,colback=teal!4,colframe=teal!35!black,title={Before you move on}]
What does das Unkraut mean? (\textquotedblleft{}Weeds\textquotedblright{}.) Say it once more.
\end{tcolorbox}
\section[jäten]{jäten — to weed}

\label{lesson:GE-C503-jaeten}



\begin{quote}
Before the new one: say the German for to garden, then the German for weeds.
\end{quote}

\subsection*{You'll want to know: jäten}
\textbf{jäten} — \textquotedblleft{}to weed\textquotedblright{}.

\begin{quote}
Am Samstag jäten wir das Unkraut im Garten.
\end{quote}

\begin{quote}
\emph{On Saturday we pull up the weeds in the garden.}
\end{quote}

Regular: \emph{hat gejätet}.

\begin{cousinweb}
\emph{jäten} is inherited from Middle High German \emph{jeten} and Old High German \emph{jetan}, from Proto-Germanic \emph{jedaną} (EtymDB).
\end{cousinweb}

\begin{grammarlens}[title={What you've built}]
One more word for telling, explaining and describing at B1.
\end{grammarlens}

\subsection*{Your turn}
\begin{itemize}
\raggedright
  \item \emph{Say it:} \emph{jäten}
  \item \emph{Say it:} \emph{jäten}, once more
  \item \emph{Say it:} a sentence of your own with \emph{jäten}
  \item \emph{Recall:} say the German for to garden, then the German for weeds, then say \emph{jäten} again
\end{itemize}

\begin{tcolorbox}[breakable,colback=teal!4,colframe=teal!35!black,title={Before you move on}]
What does jäten mean? (\textquotedblleft{}To weed\textquotedblright{}.) Say it once more.
\end{tcolorbox}
\section[die Gießkanne]{die Gießkanne — watering can}

\label{lesson:GE-C503-giesskanne}



\begin{quote}
Before the new one: say the German for weeds, then the German for to weed.
\end{quote}

\subsection*{You'll want to know: die Gießkanne}
\textbf{die Gießkanne} — \textquotedblleft{}watering can\textquotedblright{}.

\begin{quote}
Mit der Gießkanne gießt das Kind die Blumen.
\end{quote}

\begin{quote}
\emph{The child waters the flowers with the watering can.}
\end{quote}

Plural: \emph{die Gießkannen}. \emph{Gießen} — to water, to pour.

\begin{cousinweb}
\emph{Gießkanne} is built on \emph{gießen} \textquotedblleft{}to pour\textquotedblright{}, inherited from Old High German \emph{giozan}, from Proto-Germanic \emph{geutaną} (EtymDB), and \emph{Kanne} \textquotedblleft{}jug\textquotedblright{}.
\end{cousinweb}

\begin{grammarlens}[title={What you've built}]
One more word for telling, explaining and describing at B1.
\end{grammarlens}

\subsection*{Your turn}
\begin{itemize}
\raggedright
  \item \emph{Say it:} \emph{die Gießkanne}
  \item \emph{Say it:} \emph{die Gießkanne}, once more
  \item \emph{Say it:} a sentence of your own with \emph{die Gießkanne}
  \item \emph{Recall:} say the German for weeds, then the German for to weed, then say \emph{die Gießkanne} again
\end{itemize}

\begin{tcolorbox}[breakable,colback=teal!4,colframe=teal!35!black,title={Before you move on}]
What does die Gießkanne mean? (\textquotedblleft{}Watering can\textquotedblright{}.) Say it once more.
\end{tcolorbox}
\section[der Kompost]{der Kompost — compost}

\label{lesson:GE-C503-kompost}



\begin{quote}
Before the new one: say the German for to weed, then the German for watering can.
\end{quote}

\subsection*{You'll want to know: der Kompost}
\textbf{der Kompost} — \textquotedblleft{}compost\textquotedblright{}.

\begin{quote}
Die Schalen von Obst und Gemüse kommen auf den Kompost.
\end{quote}

\begin{quote}
\emph{Fruit and vegetable peel goes on the compost.}
\end{quote}

Plural, rare: \emph{die Komposte}.

\begin{cousinweb}
\emph{Kompost} was borrowed from French \emph{compost} (EtymDB), which goes back to Latin \emph{compositum} \textquotedblleft{}something put together\textquotedblright{}, from \emph{compōnere}.
\end{cousinweb}

\begin{grammarlens}[title={What you've built}]
One more word for telling, explaining and describing at B1.
\end{grammarlens}

\subsection*{Your turn}
\begin{itemize}
\raggedright
  \item \emph{Say it:} \emph{der Kompost}
  \item \emph{Say it:} \emph{der Kompost}, once more
  \item \emph{Say it:} a sentence of your own with \emph{der Kompost}
  \item \emph{Recall:} say the German for to weed, then the German for watering can, then say \emph{der Kompost} again
\end{itemize}

\begin{tcolorbox}[breakable,colback=teal!4,colframe=teal!35!black,title={Before you move on}]
What does der Kompost mean? (\textquotedblleft{}Compost\textquotedblright{}.) Say it once more.
\end{tcolorbox}
